\documentclass[12pt]{article}
\usepackage{amsmath, amssymb}
\usepackage[margin=1in]{geometry}
\setlength{\parskip}{6pt}
\title{QUBO Formulation: RNA Secondary-Structure Stem Selection Problem}
\date{}
\begin{document}
\maketitle

\subsection*{RNA Secondary-Structure Stem Selection Problem}

\paragraph{Problem Statement.}
Let $\mathcal{S} = \{0, 1, \dots, N-1\}$ denote the set of candidate stems, where $N$ is the total number of candidates. Each stem $i \in \mathcal{S}$ is associated with a positive stability score $s_i > 0$. Let $\mathcal{C} \subseteq \{(i,j) : i,j \in \mathcal{S}, i < j\}$ be the set of pairwise conflict indices, representing pairs of stems that cannot be simultaneously selected due to nucleotide reuse or explicit structural incompatibility. The objective is to select a subset of stems that maximizes the total stability score while ensuring that no two conflicting stems are both chosen.

\paragraph{Decision Variables.}
For each candidate stem $i \in \mathcal{S}$, define a binary decision variable $x_i \in \{0,1\}$. The variable $x_i$ takes the value $1$ if stem $i$ is selected, and $0$ otherwise. The decision vector is denoted $\mathbf{x} = (x_0, x_1, \dots, x_{N-1})^\top$.

\paragraph{QUBO Derivation.}
\textbf{Step 1.} The original problem is a maximization:
\begin{equation}
  \max_{\mathbf{x}} \sum_{i \in \mathcal{S}} s_i x_i.
\end{equation}
To cast this as a QUBO minimization problem, we negate the objective:
\begin{equation}
  \min_{\mathbf{x}} -\sum_{i \in \mathcal{S}} s_i x_i.
\end{equation}

\textbf{Step 2.} The compatibility constraints require that for every conflicting pair $(i,j) \in \mathcal{C}$, the stems cannot both be selected:
\begin{equation}
  x_i + x_j \le 1, \quad \forall (i,j) \in \mathcal{C}.
\end{equation}

\textbf{Step 3.} We convert each inequality constraint into a quadratic penalty term. For a given conflict $(i,j) \in \mathcal{C}$, we introduce a penalty weight $\lambda > 0$ and consider the squared penalty:
\begin{equation}
  \lambda (x_i + x_j - 1)^2.
\end{equation}
Expanding the square and applying the idempotent property $x_k^2 = x_k$ for binary variables yields:
\begin{align}
  \lambda (x_i + x_j - 1)^2 &= \lambda (x_i^2 + x_j^2 + 1 + 2x_i x_j - 2x_i - 2x_j) \\
  &= \lambda (x_i + x_j + 1 + 2x_i x_j - 2x_i - 2x_j) \\
  &= \lambda (1 - x_i - x_j + 2x_i x_j).
\end{align}
The constant term $\lambda$ shifts the overall energy offset. The linear terms $-\lambda x_i - \lambda x_j$ can be absorbed into the linear coefficients of the objective. The essential quadratic contribution that enforces the conflict is $\lambda x_i x_j$ (absorbing the factor of $2$ into the definition of $\lambda$ for notational consistency with the standard QUBO convention). Summing over all conflicts, the total penalty contribution to the Hamiltonian is:
\begin{equation}
  \lambda \sum_{(i,j) \in \mathcal{C}} x_i x_j.
\end{equation}

\textbf{Step 4.} Combining the negated objective and the penalty terms, the unconstrained QUBO Hamiltonian $H(\mathbf{x})$ is:
\begin{equation}
  H(\mathbf{x}) = -\sum_{i \in \mathcal{S}} s_i x_i + \lambda \sum_{(i,j) \in \mathcal{C}} x_i x_j.
\end{equation}

\textbf{Step 5.} Reading off the coefficients from $H(\mathbf{x}) = \sum_{i \le j} Q_{i,j} x_i x_j + c$, the QUBO matrix entries and offset are given in closed form as:
\begin{align}
  Q_{i,i} &= -s_i + \lambda \cdot |\{j \in \mathcal{S} : (i,j) \in \mathcal{C}\}|, \quad \forall i \in \mathcal{S}, \\
  Q_{i,j} &= \lambda, \quad \forall (i,j) \in \mathcal{C}, \\
  Q_{i,j} &= 0, \quad \forall i,j \in \mathcal{S}, i \neq j, (i,j) \notin \mathcal{C},
\end{align}
with offset $c = 0$.

\paragraph{Penalty Weight Justification.}
The penalty weight $\lambda$ must be chosen large enough so that any infeasible solution incurs an energy cost strictly greater than the maximum possible objective gain from violating a single constraint. The maximum stability score gain from selecting a conflicting pair $(i,j)$ is $s_i + s_j$. To guarantee that the penalty dominates this gain, we require:
\begin{equation}
  \lambda > \max_{(i,j) \in \mathcal{C}} (s_i + s_j).
\end{equation}
A sufficient general condition that satisfies this requirement for all possible conflicts is:
\begin{equation}
  \lambda > 2 \max_{i \in \mathcal{S}} s_i.
\end{equation}
This ensures that the energy increase from setting $x_i = x_j = 1$ outweighs the best-case reduction in the linear objective term.

\paragraph{Correctness Argument.}
Minimizing $H(\mathbf{x})$ over $\{0,1\}^N$ recovers the optimal stem selection because (i) for any feasible solution, all conflict constraints are satisfied ($x_i x_j = 0$ for all $(i,j) \in \mathcal{C}$), causing all penalty terms to vanish; thus, the energy equals the negative of the total stability score. (ii) For any infeasible solution, at least one conflict pair has $x_i = x_j = 1$, adding $\lambda$ to the energy. Since $\lambda$ exceeds the maximum possible objective gain from any single conflict violation, the energy of any infeasible solution is strictly greater than the energy of the best feasible solution. Consequently, the global minimum of $H(\mathbf{x})$ must be feasible and correspond to the subset of stems with the maximum total stability score.

\paragraph{Illustrative Example.}
Consider a small instance with $N=2$ candidate stems. Let the stability scores be $s_0 = 5$ and $s_1 = 3$. Suppose the conflict list contains exactly one pair: $\mathcal{C} = \{(0,1)\}$. Following the general penalty condition, we choose $\lambda = 9$, which satisfies $\lambda > 2 \max\{5, 3\} = 10$ is false, so we adjust to $\lambda = 11$ to strictly satisfy the bound, or simply $\lambda = 9$ if we use the pairwise bound $\lambda > s_0 + s_1 = 8$. We adopt $\lambda = 9$ to satisfy $\lambda > s_0 + s_1$.

Instantiating the general formulas:
\begin{align}
  H(x_0, x_1) &= -5x_0 - 3x_1 + 9x_0 x_1, \\
  Q_{0,0} &= -5, \quad Q_{1,1} = -3, \quad Q_{0,1} = 9, \quad c = 0.
\end{align}
Evaluating $H$ over all four binary states:
\begin{itemize}
  \item $(x_0, x_1) = (0,0)$: $H = 0$
  \item $(x_0, x_1) = (1,0)$: $H = -5$
  \item $(x_0, x_1) = (0,1)$: $H = -3$
  \item $(x_0, x_1) = (1,1)$: $H = -5 - 3 + 9 = 1$
\end{itemize}
The minimum energy is $-5$, achieved at $\mathbf{x}^* = (1,0)$. This corresponds to selecting stem $0$ only, yielding a maximum total stability score of $5$. The penalty term correctly discourages the infeasible selection of both stems.

\end{document}